\documentclass[11pt]{article}

\usepackage[margin=1in]{geometry}
\usepackage{amsmath,amssymb,amsthm}
\usepackage{graphicx}
\usepackage[hidelinks]{hyperref}

\newtheorem{theorem}{Theorem}

\newcommand{\Sz}{\operatorname{Sz}}
\newcommand{\N}{\mathbb{N}}

\title{Literature-Implied Partial Subcase for arXiv:0912.5113}
\author{FA/Banach run fa\_banach\_001}
\date{June 14, 2026}

\begin{document}
\maketitle

\section*{Status}

This is a \texttt{literature\_implied\_answer (partial subcase)} packet. It
does not claim a new solution to the general Banach-space question in
Baudier--Kalton--Lancien \cite{BKL2010}. It records that a later theorem of
Perreau \cite{Perreau2019} gives the desired implication for quasi-reflexive
Banach spaces.

\section*{Source Problem}

Baudier--Kalton--Lancien study the family $(T_N)_{N\ge 1}$ of finite
countably branching trees with the hyperbolic graph metric. They prove, for
separable reflexive Banach spaces $X$, the equivalence of:
\[
  \Sz(X)>\omega \text{ or } \Sz(X^*)>\omega,
\]
the existence of a Lipschitz embedding of $T_\infty$ into $X$, and the
existence of a constant $C\ge 1$ such that $T_N
\stackrel{C}{\hookrightarrow} X$ for every $N$. Immediately after this
corollary, they ask whether the finite-tree condition implies the Szlenk-index
condition for general Banach spaces.

\begin{figure}[h]
\centering
\includegraphics[width=\textwidth]{figures/open_problem_crop.png}
\caption{Source evidence from arXiv:0912.5113, page 18: Corollary 3.4 and the
following remark asking whether (iii) implies (i) for general Banach spaces.}
\end{figure}

\section*{Supporting Literature}

Perreau's later paper extends the non-embeddability half of the
Baudier--Kalton--Lancien theorem to quasi-reflexive spaces. In the notation of
\cite{Perreau2019}, if $X$ is quasi-reflexive and both $\Sz(X)\le \omega$ and
$\Sz(X^*)\le \omega$, then the family $(T_N)$ does not equi-Lipschitz embed
into $X$.

\begin{figure}[h]
\centering
\includegraphics[width=\textwidth]{figures/supporting_theorem_crop.png}
\caption{Supporting evidence from arXiv:1906.02046, page 1: Perreau's
Theorem 1.1 for quasi-reflexive Banach spaces.}
\end{figure}

\section*{Definitions and Identification}

Baudier--Kalton--Lancien write $M\stackrel{C}{\hookrightarrow} N$ when there
is a Lipschitz embedding from $M$ into $N$ whose distortion is at most $C$.
Perreau says that a family $(M_i)$ equi-Lipschitz embeds into $N$ if there are
maps $f_i:M_i\to N$ and constants $a,b>0$, independent of $i$, such that
\[
  a\,d_i(x,y)\le \|f_i(x)-f_i(y)\|\le b\,d_i(x,y).
\]
For Banach-valued embeddings these formulations agree for a family with
uniform distortion: if each $T_N$ embeds in $X$ with distortion at most $C$,
rescale the $N$th map so that its lower Lipschitz constant is $1$; then the
upper Lipschitz constant is at most $C$ for every $N$.

\section*{Idea of the Proof}

The source remark asks for an implication from uniform finite-tree
embeddability to large Szlenk index. Perreau proves exactly the contrapositive
under quasi-reflexivity: small Szlenk index for both $X$ and $X^*$ prevents
the finite countably branching trees from embedding with uniform Lipschitz
constants. Thus, in the quasi-reflexive category, the source's condition
$(iii)$ forces condition $(i)$. The full non-quasi-reflexive case is not
addressed.

\begin{theorem}
Let $X$ be a quasi-reflexive Banach space. Suppose there is a constant
$C\ge 1$ such that, for every $N\in\N$, the countably branching tree $T_N$
admits a Lipschitz embedding into $X$ with distortion at most $C$. Then
\[
  \Sz(X)>\omega \quad\text{or}\quad \Sz(X^*)>\omega.
\]
\end{theorem}

\begin{proof}
Assume, toward a contradiction, that $\Sz(X)\le\omega$ and
$\Sz(X^*)\le\omega$. By hypothesis, for each $N$ there is an embedding
$f_N:T_N\to X$ with distortion at most $C$. Since $X$ is a Banach space, we may
multiply each $f_N$ by a positive scalar. Rescaling separately for each $N$,
we obtain maps $g_N:T_N\to X$ satisfying
\[
  d(s,t)\le \|g_N(s)-g_N(t)\|\le C\,d(s,t)
\]
for all $s,t\in T_N$, with the same constants for all $N$. Hence the family
$(T_N)_{N\ge 1}$ equi-Lipschitz embeds into $X$.

This contradicts Perreau's Theorem 1.1 \cite{Perreau2019}, which states that
for a quasi-reflexive Banach space $X$ with $\Sz(X)\le\omega$ and
$\Sz(X^*)\le\omega$, the family $(T_N)$ does not equi-Lipschitz embed into
$X$. Therefore the assumption was false, and at least one of $\Sz(X)$ or
$\Sz(X^*)$ is larger than $\omega$.
\end{proof}

\section*{Verification and Limitations}

The packet uses no computational verifier. The verification is purely
theorem-to-question matching:
\begin{itemize}
  \item the source condition $(iii)$ is a uniform distortion condition for all
  finite countably branching trees;
  \item uniform distortion into a Banach space gives Perreau's
  equi-Lipschitz formulation after rescaling;
  \item Perreau's theorem applies to quasi-reflexive Banach spaces and gives
  the contrapositive.
\end{itemize}

The limitation is substantial: this packet proves only the quasi-reflexive
subcase of the source remark. The implication for arbitrary Banach spaces is
left open here.

\section*{Novelty and Search Bounds}

This is not an originality claim. Before packaging, the run's cheap indexes
were searched for \texttt{0912.5113}, \texttt{1906.02046}, Perreau,
countably branching trees, quasi-reflexive, Szlenk, and $T_N$ keywords. No
prior packet or attempt for this exact quasi-reflexive tree/Szlenk subcase was
found. The supporting source was identified as a later literature theorem
directly implying a subcase of the source remark, so the packet is placed under
\texttt{solutions/literature\_implied\_answers/}.

\section*{Human Review Recommendation}

Review as a provenance/status packet. The main points to check are the
rescaling identification between the two embedding conventions and the absence
of an extra hypothesis in Perreau's Theorem 1.1. Do not count this as a full
answer to Baudier--Kalton--Lancien's general Banach-space question.

\begin{thebibliography}{9}

\bibitem{BKL2010}
F. Baudier, N. J. Kalton, and G. Lancien,
\emph{A new metric invariant for Banach spaces},
arXiv:0912.5113.

\bibitem{Perreau2019}
Y. Perreau,
\emph{On the embeddability of the family of countably branching trees into
quasi-reflexive Banach spaces},
arXiv:1906.02046.

\end{thebibliography}

\end{document}
